\documentclass[11pt]{article}
\usepackage[utf8]{inputenc}
\usepackage[T1]{fontenc}
\usepackage[margin=1in]{geometry}
\usepackage{amsmath,amssymb}
\usepackage{booktabs}
\usepackage{listings}
\usepackage{xcolor}
\usepackage{hyperref}

\lstset{basicstyle=\small\ttfamily, breaklines=true,
        frame=single, backgroundcolor=\color{gray!8}}

\title{Additional File for Review --- Not for Publication\\
\large Geodesic Diagnostic Trajectories (GDT)}
\author{Hemanth Manchabale Papachappa \and Aniriuddha Ganguly}
\date{September 2026}

\begin{document}
\maketitle

\noindent\textit{This file contains reviewer-facing technical details that
support assessment of reproducibility and methodological soundness. It is not
intended for publication.}

\tableofcontents
\newpage

% ─────────────────────────────────────────────────────────────
\section{Computational Environment}

\begin{itemize}
  \item \textbf{Hardware:} MacBook Pro (Apple M-series), 36 GB unified memory
  \item \textbf{OS:} macOS 15.x
  \item \textbf{Python:} 3.12.4 (pyenv)
  \item \textbf{Key libraries:} PyTorch 2.x, HuggingFace Transformers,
        \texttt{rank\_bm25}, \texttt{scikit-learn}, \texttt{scipy},
        \texttt{geomstats} (SPD manifold operations)
  \item \textbf{ClinicalBERT:} \texttt{emilyalsentzer/Bio\_ClinicalBERT} from
        HuggingFace Hub
  \item \textbf{Random seeds:} All experiments seeded with $42$ for
        reproducibility
\end{itemize}

% ─────────────────────────────────────────────────────────────
\section{Data Use Agreements}

All datasets are publicly available under standard data use agreements. The
authors hold active agreements for MIMIC-III, MIMIC-IV, eICU, and HiRID via
PhysioNet. No re-identification of patient data was performed; de-identified
data were used exclusively. Synthetic vignette components generated from
aggregate distributional statistics do not constitute protected health information
under HIPAA Safe Harbor criteria.

% ─────────────────────────────────────────────────────────────
\section{Vignette Construction Protocol}

Each of the 31 test vignettes was constructed by the following procedure:
\begin{enumerate}
  \item Select a hospitalization record from MIMIC-III/IV, eICU, or HiRID with
        $\geq 3$ distinct diagnostic category transitions documented in progress
        notes or orders.
  \item Identify 3--7 ``pivot points'' in the clinical record where the
        treating team's documented focus shifted to a new problem.
  \item For each pivot point, construct a natural-language query representing
        the information need at that moment (e.g., ``current creatinine trend
        and eGFR trajectory'').
  \item Define a gold-standard target document (the note or lab report that
        best answers the query) from the de-identified record.
  \item Construct a retrieval corpus of 500--2000 documents from the same
        hospitalization plus background documents sampled from the full dataset.
\end{enumerate}

Two clinical informaticists reviewed all 31 vignettes for clinical plausibility
and query naturalness. Inter-rater agreement on gold-standard document selection
was $\kappa = 0.81$ (substantial agreement).

% ─────────────────────────────────────────────────────────────
\section{NASA-TLX Simulation Methodology}

Because no human clinicians participated, NASA-TLX subscale scores were estimated
from observable system events using the following mapping:

\begin{table}[h]
\centering
\caption{NASA-TLX subscale proxy mapping used in simulation.}
\begin{tabular}{lll}
\toprule
\textbf{Subscale} & \textbf{Observable proxy} & \textbf{Weight} \\
\midrule
Mental Demand    & Mean query complexity (token count) & 20\% \\
Temporal Demand  & Mean latency / 300s budget & 20\% \\
Effort           & Query reformulation rate & 20\% \\
Frustration      & Failed retrievals / total queries & 20\% \\
Performance      & Inverse of time-to-correct-info & 20\% \\
\bottomrule
\end{tabular}
\end{table}

These are explicitly labeled as \textit{proxy estimates}, not clinician
self-report, throughout the manuscript. We use the term ``simulated NASA-TLX''
consistently to prevent misinterpretation.

% ─────────────────────────────────────────────────────────────
\section{BM25 and TF-IDF Implementation Details}

\begin{lstlisting}[language=Python, caption=BM25 session context implementation]
from rank_bm25 import BM25Okapi
import numpy as np

class SessionBM25:
    def __init__(self, k1=1.5, b=0.75, window=10, decay=0.9):
        self.k1 = k1; self.b = b
        self.window = window; self.decay = decay

    def build_query(self, history):
        # Recency-weighted query: concatenate last `window` queries
        # with exponential downweighting of older tokens
        weighted = []
        for i, q in enumerate(history[-self.window:]):
            weight = self.decay ** (len(history[-self.window:]) - i - 1)
            weighted.extend(q.split() * max(1, int(weight * 10)))
        return weighted

    def retrieve(self, corpus, history, top_k=10):
        bm25 = BM25Okapi(corpus)
        query = self.build_query(history)
        return bm25.get_top_n(query, corpus, n=top_k)
\end{lstlisting}

% ─────────────────────────────────────────────────────────────
\section{Code Repository Structure}

Upon acceptance, the following repository structure will be made public via GitHub:

\begin{lstlisting}
gdt-clinical-retrieval/
├── src/
│   ├── models/
│   │   ├── gdt.py          # GDT retriever (SPD + GSI gate + JEPA)
│   │   ├── cma.py          # CMA comparison system
│   │   ├── bm25_session.py # Session-aware BM25 baseline
│   │   └── tfidf_session.py
│   ├── geometry/
│   │   ├── spd_manifold.py # SPD covariance, geodesic distance
│   │   └── logeuclidean.py # Log-Euclidean projection
│   ├── jepa/
│   │   ├── predictor.py    # Latent intent predictor
│   │   └── confidence.py   # Confidence scaling
│   ├── evaluation/
│   │   ├── crossover.py    # Four-arm crossover trial logic
│   │   ├── metrics.py      # Latency, NASA-TLX proxy, Prop.1 check
│   │   └── queueing.py     # M/M/c model
│   └── data/
│       ├── mimic_loader.py
│       ├── eicu_loader.py
│       └── vignette_builder.py
├── experiments/
│   ├── run_crossover.py    # Main experiment entry point
│   └── run_ablation.py
├── configs/
│   └── gdt_default.yaml    # All hyperparameters (Table S1)
├── requirements.txt
└── README.md
\end{lstlisting}

\end{document}
